% Cote 119: offprints, reports, courses, CV and letters (typed and printed part): archive cote 119, site page 62 (PDF page 63).
% offprint. Draft compared: none.
% AI-transcribed by Claude Opus 5.5 (claude-opus-5-5, adaptive thinking, effort high), Message Batches API, batch r119_full (msgbatch_01DwZNm5Y7YsKBWyAN4sFs8w), 2026-11-13; not yet checked by eye.
\scanpage{119/62}
\mspage{( 2 )}

une forme $\varphi : V \otimes V \to 1$ est appelée une forme de C-polarisation si la
forme bilinéaire $\varphi_C$ sur le $\mathbf{R}$-vectoriel V
\[ \varphi_C (x, y) = \varphi (x, C_V\, y). \]
est symétrique définie positive. L'existence d'une telle forme sur V entraîne
que V est un objet semi-simple, et la parité de $\varphi$ est $C_V^2$. Si $\varphi$ est une forme
de C-polarisation sur V, $\varphi'$ une forme de Weil sur V de même parité $C_V^2$,
$\varphi$ et $\varphi'$ sont compatibles si et seulement si $\varphi'$ est aussi une forme de C-pola-%
risation. Enfin, si $\varphi$ (resp. $\psi$) est une forme de C-polarisation sur V (resp. W),
$\varphi \oplus \psi$, $\varphi \otimes \psi$ sont des formes de C-polarisation sur $V \oplus W$, $V \otimes W$ respec-%
tivement.

Remarque 1.5. --- La définition de forme de Weil peut être généralisée
de façon évidente pour des formes bilinéaires non dégénérées $\varphi : V \otimes V \to T$
à valeurs dans un objet inversible T. Les résultats énoncés ici restent
valables dans ce contexte.

2. Polarisations.

Définition 2.1. --- Soit $\varepsilon \in Z(K)$. Une $\varepsilon$-polarisation (homogène) $\pi$
de $C$ consiste en la donnée pour chaque objet V de $C$ d'une classe d'équi-%
valence $\pi(V)$ des formes de Weil de parité $\varepsilon_V$ sur V; on exige que si
$\varphi \in \pi(V)$, $\psi \in \pi(W)$, alors
\[ \varphi \oplus \psi \in \pi (V \oplus W) \qquad \text{et} \qquad \varphi \otimes \psi \in \pi (V \otimes W). \]
L'élément $\varepsilon$ sera appelé la parité de $\pi$ et noté $\varepsilon(\pi)$.

2.2. L'ensemble des polarisations de $C$ sera noté Pol$(C)$, celui des
$\varepsilon$-polarisations, Pol$_\varepsilon(C)$. Les polarisations de parité 1 sont appelées symé-%
triques. La parité définit une application
\[ (2.2.1) \qquad \varepsilon : \ \mathrm{Pol}(C) \to Z(\mathbf{R}). \]
Si Pol$(C) \neq \emptyset$ on prouve facilement que la catégorie $C$ est semi-simple.
Si $C$ est présentée comme réunion filtrante de ses sous-catégories tanna-%
kiennes algébriques [voir (${}^1$), 2.3], $C = \varinjlim C_i$, on prouve qu'on a
\[ (2.2.2) \qquad \mathrm{Pol}(C) \simeq \varprojlim \mathrm{Pol}(C_i). \]
D'autre part, si $\varepsilon \in Z(K)$, il résulte de (1.3.1) qu'on a une bijection cano-%
nique
\[ (2.2.3) \qquad \mathrm{Pol}_\varepsilon(C) \simeq \mathrm{Pol}_\varepsilon(C_{\mathbf{R}}), \] \margin{\unsure{/}}
où $C_{\mathbf{R}}$ est la catégorie tannakienne sur $\mathbf{R}$ obtenue par extension des
scalaires de K à $\mathbf{R}$ (${}^2$). Les isomorphismes (2.2.2), (2.2.3) justifient la
convention suivante :

Convention 2.3. --- Dans la suite de cette Note, $C$ est une catégorie tanna-%
kienne algébrique sur $\mathbf{R}$.

\mspage{\hand{62}}
